\SourceLeaf{419}%
\noindent so gilt
\[
(P\cos\varphi^2+2Q\cos\varphi\sin\varphi+R\sin\varphi^2)ds=(p\sin\varphi-q\cos\varphi)d\varphi,
\]
also
\[
\frac1r=\frac{P\cos\varphi^2+2Q\cos\varphi\sin\varphi+R\sin\varphi^2}{p\sin\varphi-q\cos\varphi}
\]
oder, da
\[
\cos\varphi=\frac{\mp q}{\BareRoot{(pp+qq)}},\qquad\sin\varphi=\frac{\pm p}{\BareRoot{(pp+qq)}}:
\]
\[
\pm\frac1r=\frac{Pqq-2Qpq+Rpp}{(pp+qq)^{\frac32}}.
\]

\Needspace{5\baselineskip}
\begin{center}4.\end{center}
\noindent\hspace*{1.2em}Das zweideutige Vorzeichen in der letzten Formel mag zunächst auffallen, erweist sich aber bei näherer Betrachtung als völlig in Ordnung. Dieser Ausdruck hängt nämlich nur von den partiellen Differentialen von $V$ ab, und die Funktion $V$ stellt für sich nur die Gestalt der gekrümmten Linie dar, ohne anzugeben, in welchem Sinn man sie sich als durchlaufen denkt. Daher muss die Frage, ob die Kurve nach rechts oder nach links konvex ist, offenbleiben, bis dieser Sinn auf andere Weise bestimmt ist. Ebenso verhält es sich mit der Bestimmung von $\varphi$ durch seine Tangente, bei der für sich genommen eine Unbestimmtheit von $180^{0}$ bleibt. Der Sinn, in dem die Kurve durchlaufen wird, lässt sich auf verschiedene Arten angeben.

I. Durch das Vorzeichen der Änderung von $x$. Nimmt $x$ zu, so muss $\cos\varphi$ positiv sein. Daher gelten die oberen Vorzeichen, wenn $q$ einen negativen Wert hat, und die unteren, wenn $q$ einen positiven Wert hat. Nimmt $x$ ab, gilt das Umgekehrte.

II. Durch das Vorzeichen der Änderung von $y$. Nimmt $y$ zu, so erfordert ein positiver Wert von $p$ die oberen und ein negativer die unteren Vorzeichen. Umgekehrt verhält es sich bei abnehmendem $y$.

III. Durch das Vorzeichen des Werts, den die Funktion $V$ annimmt, wenn man die Kurve verlässt. Die Änderungen von $x,y$ beim Verlassen der Kurve senkrecht zur Tangente nach rechts, also in Richtung $\varphi+90^{0}$, seien $\delta x,\delta y$, bei der Länge der Normalen $=\delta\varrho$. Man hat dann\space%
